\documentclass[11pt]{article}
\usepackage[T1]{fontenc}
\usepackage{lmodern,amsmath,amssymb,amsthm,array,booktabs}
\usepackage[margin=1in]{geometry}
\usepackage[hidelinks,breaklinks]{hyperref}
\newtheorem{theorem}{Theorem}
\newtheorem{corollary}[theorem]{Corollary}
\newcommand{\id}{\mathrm{id}}
\newcommand{\spec}{\operatorname{spec}}
\newcommand{\tr}{\operatorname{Tr}}
\newcommand{\E}{\mathcal E}
\newcommand{\A}{\mathcal A}
\newcommand{\ket}[1]{\lvert #1\rangle}
\newcommand{\bra}[1]{\langle #1\rvert}
\title{Quantitative extensions of Ara\'ujo et al.'s phase obstruction
for Liu, Weso\l{}owski, Thompson, Gu and Pira's
uncontrolled fractional-evolution question}
\author{Haobo Ma\\
\small ChronoAI Pte Ltd, Singapore
\and Wenlin Zhang\\
\small National University of Singapore\\
\small ChronoAI Pte Ltd, Singapore}
\date{16 November 2026\\[3pt]
\small Creative Commons Attribution 4.0 International (CC BY 4.0)}
\begin{document}
\maketitle
\begin{abstract}
Extending Ara\'ujo, Feix, Costa and Brukner's global-phase argument
(2014), we answer negatively the question in Liu, Weso\l{}owski, Thompson,
Gu and Pira, \emph{Optimal query complexity for fractional quantum evolution},
arXiv:2610.01940v1 (1 October 2026), Section V, page 7, unnumbered
``Uncontrolled access'' paragraph, with its prescribed branch and uniform
channel criterion.
For $0<t<1$, $0<\delta<\pi$ and every dimension $d\geq2$, ordinary
queries have an unhalved diamond-error floor $\sin(\pi t)$, independent
of their number. Retaining a pure full output strengthens the bound to
$2\sin(\frac\pi2\min\{t,1-t\})$.
The same argument gives the endpoint order of the best attainable error
and classifies continuously differentiable phase functions on this branch:
arbitrarily accurate ordinary-access implementation is possible exactly
for integer powers, up to a constant phase.
\end{abstract}

\section{Question, prior results and conventions}
Controlled access makes a global phase measurable as a relative phase;
ordinary access does not. Ara\'ujo et al.~\cite[p.~2]{araujo} use this
distinction to obstruct universal control of an unknown operation.
Sheridan, Maslov and Mosca~\cite[Section II.A and Appendix A]{sheridan}
develop fractional evolution with controlled queries and explain the
different, eigenphase-relative target available from a reference register.
Gavorov\'a, Seidel and Touati~\cite[Section IV.C, Corollary 1]{gavorova}
prove an all-domain root impossibility by topological methods.
These results supply the phase-obstruction principle and its broader context.
The point here is an explicit admissible pair on the promised branch,
its quantitative separation, and consequences for other phase functions.

In Section V of~\cite{liu}, the source states:
``It is unclear whether the matching upper bound survives without controlled
access.'' The matching bound is $O(\delta^{-1}\log(1/\varepsilon))$.
This is a question, not a claimed ordinary-access theorem. We answer it with
Definitions 1--3 and Appendix B of that source, replacing controlled slots
by ordinary slots. Its controlled-access results are separate.

Let $S=\mathbb C^d$, let $A=\mathbb C^r$ with $1\leq r<\infty$,
and fix a unit vector $\ket0_A$. Write
$J\ket\psi=\ket0_A\otimes\ket\psi$.
For Hermitian $H$ with $\spec(H)\subset[\delta,2\pi)$, put
$U=e^{iH}$ and $U^t=e^{itH}$. The half-open branch uniquely fixes $H$.
A fixed ordinary-query circuit has the form
\begin{equation}\label{eq:circuit}
 \A_U=W_N(I_A\otimes U^{\sigma_N})W_{N-1}\cdots
 W_1(I_A\otimes U^{\sigma_1})W_0,
 \qquad \sigma_j\in\{-1,+1\},
\end{equation}
where the $W_j$ are unitary and independent of $U$; $N=0$ is allowed.
Negative signs mean adjoint queries. Gates may depend on $t,\delta$ and
the requested accuracy.
For an isometry $V$, put $\E_V(X)=VXV^\dagger$. Our diamond norm is
unhalved: for differences of channels from $S$, with $R=\mathbb C^d$,
\[
 \|\Phi\|_\diamond
 =\sup_{\rho\geq0,\,\tr\rho=1}
   \|(\Phi\otimes\id_R)(\rho)\|_1,
 \qquad \rho\text{ on }S\otimes R.
\]
The error of a circuit is
\begin{equation}\label{eq:error}
 e(\A;t,\delta)=
 \sup_{\substack{H=H^\dagger\\\spec(H)\subset[\delta,2\pi)}}
 \|\E_{\A_{e^{iH}}J}-\E_{Je^{itH}}\|_\diamond.
\end{equation}
Thus the comparison includes the initialized workspace, as in Appendix B
of~\cite{liu}; it does not merely compare a reduced system output.

\section{A promised phase pair}
\begin{theorem}\label{thm:floor}
Fix $0<t<1$, $0<\delta<\pi$ and $d\geq2$.
Every circuit~\eqref{eq:circuit} satisfies
\[
 e(\A;t,\delta)\geq\sin(\pi t).
\]
The same lower bound holds for a deterministic algorithm built from ordinary
oracle channels, oracle-independent quantum operations and classical
conditioning, when its system output is compared to $\E_{e^{itH}}$.
No oracle-dependent auxiliary input is supplied in either model.
\end{theorem}
\begin{proof}
First take $d=2$ and use
\[
 H_+=\operatorname{diag}(\delta,\delta+\pi),\qquad
 H_-=\operatorname{diag}(\delta+\pi,\delta).
\]
Both spectra lie in $[\delta,2\pi)$ because $\delta+\pi<2\pi$.
Their exponentials obey $U_-=-U_+$.
For either sign, $(-U_+)^{\sigma_j}=-U_+^{\sigma_j}$.
Induction on the ordered product in~\eqref{eq:circuit} therefore gives
\begin{equation}\label{eq:phase}
 \A_{U_-}=(-1)^N\A_{U_+},\qquad
 \E_{\A_{U_-}J}=\E_{\A_{U_+}J}.
\end{equation}
The base case is $\A_U=W_0$; each next slot contributes precisely one
additional scalar $-1$, which commutes through its fixed gate.

For $\ket+=(\ket0+\ket1)/\sqrt2$, let $\rho_+=\ket+\bra+$.
The ideal system outputs are
\[
 R_t=\frac12\begin{pmatrix}1&e^{-i\pi t}\\e^{i\pi t}&1\end{pmatrix},
 \qquad
 S_t=\frac12\begin{pmatrix}1&e^{i\pi t}\\e^{-i\pi t}&1\end{pmatrix}.
\]
Hence
\[
 D_t=R_t-S_t
 =\begin{pmatrix}0&-i\sin(\pi t)\\i\sin(\pi t)&0\end{pmatrix},
 \quad D_t^\dagger D_t=\sin^2(\pi t)I_2,
 \quad \|D_t\|_1=2\sin(\pi t).
\]
Isometric embedding by $J$ and tensoring a fixed reference state preserve
trace norm. On this input let $X$ be the common implemented output from
\eqref{eq:phase}, and let $e_+$ and $e_-$ be the two diamond errors. Then
\[
 2\sin(\pi t)\leq
 \|JR_tJ^\dagger-X\|_1+\|X-JS_tJ^\dagger\|_1
 \leq e_++e_-.
\]
Here the input can be tensored with a fixed reference state in the diamond
supremum. Thus one error is at least $\sin(\pi t)$.

For $d>2$, append the entry $\delta$ to $H_+$ and $\delta+\pi$ to $H_-$
in every extra coordinate. The sign relation persists, and the same
input is supported on the first two coordinates.

For the more general model, the oracle channels for $U_+$ and $U_-$
coincide, as do those for their adjoints. Substituting equal channels
at every slot leaves all outcome probabilities and conditional operations
unchanged. Thus the final deterministic output channels coincide, even
with discarded workspace or classical randomness. The same triangle
inequality on the system output proves the claim.
\end{proof}

\begin{corollary}\label{cor:pure}
In the full-output unitary-circuit model~\eqref{eq:circuit}, put
$s=\min\{t,1-t\}$. Under Theorem~\ref{thm:floor}'s hypotheses,
\[
 e(\A;t,\delta)\geq \sqrt{2-2|\cos(\pi t)|}
 =2\sin(\pi s/2).
\]
In particular the square-root error is at least $\sqrt2$.
\end{corollary}
\begin{proof}
The two ideal full outputs on $\ket+$ are unit vectors $p,q$ with
$|\langle p,q\rangle|=|\cos(\pi t)|=:c$.
The common implemented state is also pure, say $\ket x\bra x$.
The nonzero eigenvalues of $\ket p\bra p+\ket q\bra q$ are $1+c$ and
$1-c$, as follows from the Gram matrix of $p,q$.
Therefore
$|\langle x,p\rangle|^2+|\langle x,q\rangle|^2\leq1+c$,
so one squared overlap is at most $(1+c)/2$.
For unit vectors $v,w$, the possible nonzero eigenvalues of
$\ket v\bra v-\ket w\bra w$ are
$\pm\sqrt{1-|\langle v,w\rangle|^2}$.
The corresponding trace-norm error is consequently at least
$2\sqrt{(1-c)/2}$, and a diamond error is no smaller.
Finally $|\cos(\pi t)|=\cos(\pi s)$ for $0<t<1$.
\end{proof}

\begin{corollary}\label{cor:source}
The ordinary-access question in Section V of~\cite{liu} has answer no:
there is no arbitrarily accurate uniform implementation, and hence no
matching uniform query upper bound. At $d=2$, $t=\delta=1/2$ and
$\varepsilon=1/4$, no finite ordinary-query circuit succeeds.
\end{corollary}
\begin{proof}
The source's range $0<\delta<1$ is contained in $0<\delta<\pi$.
At $t=1/2$, Theorem~\ref{thm:floor} already gives error at least $1>1/4$,
independently of query count. The stronger bound in Corollary~\ref{cor:pure}
uses the full-output purity in addition.
\end{proof}

\section{Endpoint dependence and other phase functions}
\begin{corollary}\label{cor:endpoint}
Let $E_*(t,\delta,d)$ be the infimum of~\eqref{eq:error} over all finite
circuits, all finite workspaces and all query counts. If
$d\geq2$, $0<\delta<\pi$ and $0<t<1$, then
\[
 2\sin(\pi s/2)\leq E_*(t,\delta,d)
 \leq\min\{2,(2\pi-\delta)s\},\qquad s=\min\{t,1-t\}.
\]
Thus $E_*=\Theta(s)$ with absolute constants. At $t=0$ and $t=1$ the
error is zero, using respectively zero queries and one ordinary query.
\end{corollary}
\begin{proof}
The lower bound is Corollary~\ref{cor:pure}. For the upper bound put
$b=(\delta+2\pi)/2$. Every promised $H$ obeys
$\|H-bI\|_{\rm op}\leq(2\pi-\delta)/2$.
For $t\leq1/2$, the identity circuit approximates $e^{itH}$ after a
scalar phase with operator-norm error at most
$t\|H-bI\|_{\rm op}$, by $|e^{iu}-1|\leq|u|$ on each eigenvalue.
For $t\geq1/2$, the one-query circuit $U$ gives the same estimate with
$1-t$. For two isometries $V,W$, adding and subtracting $W\rho V^\dagger$
in their channel difference gives
$\|\E_V-\E_W\|_\diamond\leq2\|V-W\|_{\rm op}$;
the bound also holds after a reference tensor factor. Channel distance
is at most $2$. Finally $2s\leq2\sin(\pi s/2)\leq\pi s$ for
$0\leq s\leq1/2$, which proves the stated order.
\end{proof}
The upper-bound observation is already stated in~\cite[Section V,
``Dependence on $t$'']{liu}. The corollary combines it with the ordinary-access
lower bound. It concerns the best error, not the controlled-access
query-complexity conjecture in that paragraph.

\begin{corollary}\label{cor:functions}
Fix $d\geq2$ and $0<\delta<\pi$. Let
$f:[\delta,2\pi)\to\mathbb R$ be continuously differentiable, with the
one-sided interpretation at $\delta$. The target channel
$\E_{e^{if(H)}}$ admits arbitrarily accurate uniform approximation from
ordinary queries, without an oracle-dependent auxiliary input, if and only if
\[
 f(\theta)=m\theta+b\quad(\delta\leq\theta<2\pi)
 \quad\text{for some }m\in\mathbb Z,\ b\in\mathbb R.
\]
For these functions, $|m|$ signed queries implement the target exactly.
\end{corollary}
\begin{proof}
Any two admissible oracles differing by a scalar phase supply the same
ordinary query channels. If arbitrarily small uniform errors are possible,
their target channels must coincide: their diamond separation is bounded
by twice every such error.

Choose any $x,y\in(\delta,2\pi)$. For every sufficiently small real $h$,
both $H=\operatorname{diag}(x,y)$ and $H+hI$ are admissible, and their
oracles differ by $e^{ih}$. Applying equality of target channels to the
off-diagonal matrix unit gives
\[
 e^{i(f(x+h)-f(y+h))}=e^{i(f(x)-f(y))}.
\]
Differentiating at $h=0$ yields $f'(x)=f'(y)$.
Consequently $f(\theta)=a\theta+b$ on the open interval, and continuity
extends this to $\delta$.
Apply target-channel equality also to the sign pair in
Theorem~\ref{thm:floor}. Its off-diagonal entries are $e^{-ia\pi}$ and
$e^{ia\pi}$, so $e^{2\pi ia}=1$ and $a\in\mathbb Z$.
For $d>2$ append repeated admissible eigenvalues to these two-dimensional
constructions; the same matrix unit gives both conclusions.
Conversely, $e^{if(H)}=e^{ib}U^m$ induces the channel of $U^m$.
Use $m$ forward queries when $m>0$, $-m$ adjoint queries when $m<0$,
and no queries when $m=0$.
\end{proof}
This ordinary-access classification is not asserted in the posing source.
It supplies an obstruction for analytic phase functions as a special case,
but does not determine their controlled-access complexity.

\paragraph{Resources and boundary cases.}
The strict endpoint $\delta<\pi$ matters to the displayed sign pair:
at $\delta=\pi$ its other eigenvalue is $2\pi$, outside the promise.
No impossibility for $\delta\geq\pi$ is inferred here.
For $d=1$ all target channels are the identity, so the restriction $d\geq2$
is necessary. An oracle-dependent reference register can change the target:
Sheridan et al.'s construction uses an eigenphase-relative power, not this
uniform prescribed-branch map~\cite[Appendix A]{sheridan}.
Controlled queries also change the information supplied: the controlled
versions of $U$ and $-U$ are not equal channels. For the fixed pair
$\operatorname{diag}(i,-i)$ and its negative, a superposed control and
the first basis vector yield orthogonal outputs. Thus the ordinary-channel
argument leaves the controlled results intact.

\section{Source coverage and mathematical verification}
All source locators in the table refer to~\cite{liu}; Section V is on page 7.
\begin{center}\small
\begin{tabular}{@{}>{\raggedright\arraybackslash}p{0.32\linewidth}>{\raggedright\arraybackslash}p{0.62\linewidth}@{}}
\toprule
Source item & Coverage\\
\midrule
Uncontrolled-access question & Answered: no. Theorem~\ref{thm:floor} and
Corollary~\ref{cor:source}; all source gaps and all finite query counts.\\[3pt]
Endpoint-error observation in ``Dependence on $t$'' & Reproved; prior
observation Liu et al., 2026. Corollary~\ref{cor:endpoint} adds a lower bound
of the same order for ordinary access.\\[3pt]
Controlled adaptive-in-$t$ conjecture & Not covered. An error floor is
not a controlled query-complexity estimate.\\[3pt]
``Other spectral transformations'' & The controlled analytic-function
complexity conjecture is not covered. The ordinary-access classification
in Corollary~\ref{cor:functions} is not posed; proved here.\\[3pt]
``Beyond the worst case'' & Not covered. Neither an optimal Haar-average
complexity nor a typical-oracle error follows from this pair.\\[3pt]
Theorems 1--4; Appendix C & Controlled-access results not reassessed;
the source's lower-bound implication for ordinary access is compatible
with an error floor.\\
\bottomrule
\end{tabular}
\end{center}

The global-phase step is reproved from the principle of
Ara\'ujo et al.~\cite{araujo}; the promised pair and its separation make
the present application explicit.
Theorem~\ref{thm:floor} and Corollaries~\ref{cor:pure}--\ref{cor:functions}
have written proofs above.
The pinned Lean declaration
\begin{center}\footnotesize
\texttt{Trureturing.OpenProblems.FractionalEvolutionPhase.result}
\end{center}
states the negation of the uniform ordinary-query upper-bound assertion,
using $d=2$, $t=\delta=1/2$, $\varepsilon=1/4$, arbitrary finite nonzero
workspace and the full-output unhalved criterion.
Its concrete logarithms are
$\operatorname{diag}(\pi/2,3\pi/2)$ and the swapped diagonal.
It does not state the parameterized error floor, the stronger pure-output
bound, the enlarged gap range or the phase-function classification.
\href{https://github.com/the-omega-institute/trureturing/blob/42aeb4e95344fb1ea9d2042d02ff1a5589faaacb/D5/S3/Quantum/QuantumChannels/FractionalEvolutionPhaseObstruction.lean}{Pinned Lean source for this declaration}.
The exact-arithmetic checker supplies finite matrix identities for the
square-root witness and ordered circuits. Those finite checks supplement,
rather than replace, the universal written arguments.

\begin{thebibliography}{9}
\bibitem{liu}
A. Yuezhang Liu, A. Weso\l{}owski, J. Thompson, M. Gu and L. Pira,
\emph{Optimal query complexity for fractional quantum evolution},
\href{https://arxiv.org/abs/2610.01940v1}{arXiv:2610.01940v1} (2026).
Definitions 1--3, Section V (p.~7) and Appendix B.
\bibitem{araujo}
M. Ara\'ujo, A. Feix, F. Costa and \v C. Brukner,
\emph{Quantum circuits cannot control unknown operations},
New Journal of Physics \textbf{16}, 093026 (2014),
\href{https://doi.org/10.1088/1367-2630/16/9/093026}{doi:10.1088/1367-2630/16/9/093026}.
Preprint: \href{https://arxiv.org/abs/1309.7976v2}{arXiv:1309.7976v2}.
\bibitem{sheridan}
L. Sheridan, D. Maslov and M. Mosca,
\emph{Approximating fractional time quantum evolution},
Journal of Physics A: Mathematical and Theoretical \textbf{42}, 185302 (2009),
\href{https://doi.org/10.1088/1751-8113/42/18/185302}{doi:10.1088/1751-8113/42/18/185302}.
\href{https://arxiv.org/abs/0810.3843v2}{arXiv:0810.3843v2},
Section II.A and Appendix A.
\bibitem{gavorova}
Z. Gavorov\'a, M. Seidel and Y. Touati,
\emph{Topological obstructions to quantum computation with unitary oracles},
Physical Review A \textbf{109}, 032625 (2024),
\href{https://doi.org/10.1103/PhysRevA.109.032625}{doi:10.1103/PhysRevA.109.032625}.
\href{https://arxiv.org/abs/2011.10031v3}{arXiv:2011.10031v3},
Section IV.C, Corollary 1.
\end{thebibliography}
\end{document}
